%% ============================================================================
\section{Reproduction of the Reference Configuration}
\label{app:reproduction}
\addcontentsline{lot}{section}{Reproduction of the Reference Configuration}
%% ============================================================================

This appendix documents an independent reproduction of the reference configuration
of Casini et al.~\cite{Casini23}, undertaken on a single-GPU workstation before the
cluster experiments.
It is included for two reasons.
First, the accuracy floor asserted in Section~\ref{sec:dataset} is only meaningful if
it has been established rather than assumed.
Second, the reproduction surfaced a set of pipeline defects
(Appendix~\ref{app:defects}) and one instructive environmental sensitivity, both of
which bear on how the primary results should be read.

\subsection{Environment}

Reproduction was performed on an NVIDIA RTX~3050 Laptop GPU (4~GB VRAM,
driver 566.07), 8 logical CPU cores and 15.7~GB system RAM, running
PyTorch 2.13.0+cu126 with CUDA 12.6 and Lightning 2.6.5 under Python 3.12.
The corpus was the Bing 1k configuration: 4{,}734 site tiles and 1{,}155 empty
tiles, 5{,}889 image/mask pairs totalling 957~MB, split 4{,}712 train /
588 validation / 589 test under seed 1234.

The site count is consistent with the unfiltered variant of the published corpus:
$4{,}934 - 200 = 4{,}734$, i.e. the top-200-by-area removal applied but not the
684-site minimum-size filter.
The applicable reference figure is therefore Model~1 of
Table~\ref{tab:casini-reference}, $74.17 \pm 0.38$.

\subsection{Faithfulness of the port}
\label{sec:repro-faithfulness}

The reference implementation was ported directly rather than reimplemented.
The following were held identical to the published notebook: MA-Net with an
EfficientNet-B3 ImageNet-pretrained encoder; \code{FocalLoss(mode=BINARY)} with
$\alpha$ unset; \code{Adam} at $10^{-4}$; \code{precision=32}; 20 epochs; a
\code{RandomCrop(512)} followed by \code{Resize(256)} for both training and
validation; augmentation by single-axis flip, 90-degree rotation and
brightness/contrast jitter, all at $p = 0.25$; standardisation of raw
0--255 input by ImageNet statistics without prior division by 255; evaluation by
\code{smp.metrics} with \code{reduction="macro-imagewise"}; final-epoch weights
rather than best-validation weights; and ten test passes with independent random
crops.

Two adaptations were unavoidable.
The published notebook targets PyTorch Lightning 1.x, whose
\code{training\_epoch\_end} and \code{validation\_epoch\_end} hooks were removed in
2.x; per-epoch aggregation was therefore reimplemented explicitly.
The \code{A.Flip} transform has since been removed from \code{albumentations};
its semantics - selecting one of \{vertical, horizontal, both\} uniformly, rather
than applying two independent flips - were reproduced explicitly and verified
empirically over 4{,}000 samples, yielding 0.752 no-flip and 0.089 / 0.084 / 0.075
for the three flip modes against an expected 0.750 / 0.083 each.

\subsection{Result and the batch-size constraint}
\label{sec:repro-result}

Under this configuration the reproduction attained an IoU of
$\mathbf{0.4470 \pm 0.0059}$ under the reference reduction, against the published
$0.7417 \pm 0.0038$.

The discrepancy is attributable to accelerator memory.
The reference configuration trains at batch size 32.
Measured peak allocation for MA-Net/EfficientNet-B3 at $256 \times 256$ in fp32 is
2.44~GB at batch 8, 4.65~GB at batch 16, and out-of-memory at batch 32 on a 4~GB
card; at half precision, batch 32 requires 4.94~GB and therefore also does not fit.
Batch 32 is thus unreachable on this hardware at any precision.
The reproduction used batch 8 with gradient accumulation $\times 4$, which preserves
the effective batch - and hence the gradient signal - but leaves batch
normalisation estimating its statistics from 8 samples per forward pass rather
than 32.

EfficientNet is batch-normalisation dense, and the degradation of BN statistics at
small batch is a well-characterised effect; it is the motivating problem for
Group Normalization~\cite{Wu18,Ioffe15}.
That this is the operative cause rather than a defect in the port is corroborated by
the cluster measurements of Section~\ref{sec:measured-baseline}, where the same code
at batch 32 on a V100 attains 68--73\% IoU, consistent with the published figure.

\paragraph{Consequence for this thesis.}
All accuracy figures reported in the main text are drawn from the cluster.
The workstation environment is used for development, defect analysis and
instrumentation work, for which it is adequate, and not for any claim about
segmentation quality.

\subsection{Closing the gap: full-protocol replication on the cluster}
\label{sec:repro-closure}

The batch-size hypothesis of Section~\ref{sec:repro-result} was tested directly
rather than left as an inference from BatchNorm theory. A from-scratch replica,
\code{0-reproduction}, was built independently of both the workstation port
above and the main \pads{} pipeline evaluated in Sections~\ref{ch:results}
and~\ref{ch:ablation}, for the specific purpose of isolating implementation
differences from a genuine reproducibility gap. It was run on the cluster at the
reference's full protocol: batch size 32, 20 epochs with no early stopping,
final-epoch weights, and ten independent test-set passes -- for all three
architectures used elsewhere in this thesis, not MA-Net alone.

\begin{table}[!ht]
\centering
\begin{tabular}{lcc}
\toprule
Architecture & Test IoU (mean, 10 passes) & Range \\
\midrule
DeepLabV3+ & 0.728 & 0.722--0.732 \\
SegFormer  & 0.733 & 0.725--0.739 \\
MA-Net     & 0.735 & 0.729--0.742 \\
\midrule
Reference (Casini et al., Model~1) & \textbf{0.7417} & $\pm 0.0038$ \\
\bottomrule
\end{tabular}
\caption{Full-protocol replica results, all three architectures, against the
reference figure. Standard deviation across the ten test passes was 0.003--0.004
for each architecture -- tight enough that the result is not an artefact of a
favourable single run.}
\label{tab:repro-full}
\end{table}

All three architectures land within 1--1.4 percentage points of the reference
figure. This closes the gap the workstation reproduction could not: the earlier
result (0.4470, then 0.4504 under the reference's own metric convention -- see
Section~\ref{sec:repro-metric}) was constrained by batch size 8, and the
diagnosis in Section~\ref{sec:repro-result} is confirmed directly rather than
only by analogy to published BatchNorm sensitivity results. It also settles a
question that diagnosis alone could not: whether a correct batch size on its own
would be sufficient, or whether some other, unidentified defect was also
contributing. It was sufficient -- once batch size, the full epoch budget, and
final-epoch (rather than early-stopped) weights were all restored together, no
further correction was needed to reach reference-comparable accuracy on any of
the three architectures.

This result is the basis for the caveat in Section~\ref{sec:res-accuracy}: the
main \pads{} pipeline's IoU figures throughout Sections~\ref{ch:results}
and~\ref{ch:ablation} (clustering around 0.32--0.40) are not this workload's
ceiling. They reflect a pipeline that, independently of the batch-size question
addressed here, differs from both the reference and this appendix's own replica
in two further respects: an \code{IoU} averaging convention that pools
confusion-matrix counts across the validation split rather than averaging
per-image (Section~\ref{sec:repro-metric} quantifies a 3.7-point effect from this
alone on one checkpoint), and a flip-augmentation implementation that applies two
independent single-axis flips rather than the reference's \code{A.Flip}
semantics (Section~\ref{sec:defect-augmentation}, Appendix~\ref{app:defects}).
Both are properties of \code{3-pads/*/model.py} and \code{3-pads/*/dataset.py}
specifically, not of this appendix's already-faithful port or of
\code{0-reproduction}; neither pipeline used to produce
Table~\ref{tab:repro-full} exhibits either defect. The two effects are not
individually large enough to explain the full gap between $\sim$0.35 and
$0.73$--$0.74$ on their own, and combined with an abbreviated training protocol
(early stopping rather than a fixed 20 epochs, in the main pipeline's evaluation
runs) they account for the remainder.

\subsection{Metric reduction}
\label{sec:repro-metric}

The reduction convention accounts for a material part of the difference between
naively reported and reference-comparable numbers.
Measured on the same trained checkpoint (Table~\ref{tab:reduction}):

\begin{table}[!ht]
\centering
\begin{tabular}{lc}
\toprule
Reduction & IoU \\
\midrule
Globally pooled, site-bearing tiles only          & 0.4131 \\
Per-image mean, site-bearing tiles only           & 0.3364 \\
Per-image mean, all tiles (reference convention)  & \textbf{0.4504} \\
\bottomrule
\end{tabular}
\caption{Effect of the aggregation convention on reported IoU for an identical
model. The reference convention counts a correctly-predicted empty tile as a
perfect match.}
\label{tab:reduction}
\end{table}

Two hypotheses about the metric were tested and rejected.
Excluding the empty tiles entirely changes pooled IoU by $0.001$, because they
contribute only 0.8\% of false positives - the model classifies 102 of 104 empty
tiles perfectly.
Sweeping the decision threshold over $\{0.2, \dots, 0.7\}$ confirms $0.5$ is optimal,
so the reported figure is not an artefact of a poorly chosen operating point.

\subsection{Error structure}

The aggregate figure conceals a strongly bimodal error distribution.
On the validation split, per-image IoU has a 90th percentile of $0.785$ - at
reference level - while 35.3\% of tiles score below $0.1$.
Performance is strongly monotonic in the fraction of the crop occupied by the site
(Table~\ref{tab:by-area}).

\begin{table}[!ht]
\centering
\begin{tabular}{lccc}
\toprule
Site area in crop & Tiles & Mean IoU & Fraction $< 0.1$ \\
\midrule
0.5--2\%    &  25 & 0.024 & 0.92 \\
2--5\%      &  83 & 0.179 & 0.57 \\
5--12\%     & 137 & 0.285 & 0.37 \\
12--30\%    & 151 & 0.440 & 0.19 \\
30--100\%   &  86 & 0.488 & 0.23 \\
\bottomrule
\end{tabular}
\caption{Per-image IoU by the fraction of the cropped tile occupied by the
annotated site.}
\label{tab:by-area}
\end{table}

This is consistent with the original authors' rationale for excluding sites below
1{,}000~m\textsuperscript{2}: on small objects, a few pixels of boundary error
destroys the ratio, and such annotations frequently denote sites of imprecise
location rather than precise small mounds.
It also explains the $+3.9$ point gain the reference reports from filtering
(Model~1 $\rightarrow$ Model~2).
